\documentclass{article}
\usepackage[UTF8]{ctex}
\usepackage{amsmath,amssymb,amsfonts}
\usepackage{enumitem}
\usepackage{geometry}
\geometry{a4paper, margin=1in}
% === 错题本排版优化（不参与题目标识） ===
\setlength{\parindent}{0pt}
\setlength{\parskip}{0.35em plus 0.10em minus 0.08em}
\setlength{\emergencystretch}{3em}
\setlist{itemsep=0.25em, topsep=0.40em, parsep=0pt}
\allowdisplaybreaks[2]
% === 错题本排版优化结束 ===
\begin{document}
\section*{西南财经大学2024年高等代数真题}
\section*{164．西南财经大学2024年高等代数试卷（回忆版）}
\begin{itemize}

\subsection{题目1}
  \item[1．] 求行列式
\[
D=\left|\begin{array}{ccccccc}
1 & 2 & 4 & \cdots & 2^{n-3} & 2^{n-2} & 2^{n-1} \\
0 & 0 & 0 & \cdots & 0 & 1 & -2 \\
0 & 0 & 0 & \cdots & 1 & -2 & 0 \\
\vdots & \vdots & \vdots & & \vdots & \vdots & \vdots \\
0 & 0 & 1 & \cdots & 0 & 0 & 0 \\
0 & 1 & -2 & \cdots & 0 & 0 & 0 \\
1 & -2 & 0 & \cdots & 0 & 0 & 0
\end{array}\right|
\]

\subsection{题目2}
  \item[2．] 已知矩阵 \(A\) 的秩 \(r(A)=r\) ，证明：存在 \(B_{1}, B_{2}, \cdots, B_{r}\) ，满足 \(r\left(B_{i}\right)=1(i=1,2, \cdots, r)\) ，且
\[
A=B_{1}+B_{2}+\cdots+B_{r} .
\]
\end{itemize}

\subsection{题目3}
3．求非齐次线性方程组

\[
\left\{\begin{array}{l}
x_{1}+x_{2}+x_{3}=1 ; \\
x_{1}+? x_{2}+a x_{3}=1 ; \quad(\text { 具体数据记不清楚 }) \\
2 x_{1}+? x_{2}+? x_{3}=b .
\end{array}\right.
\]

讨论 \(a, b\) 取何值时，方程组无解，有唯一解，有无穷多解，并在有无穷多解时求通解．

\subsection{题目4}
4．设 \(\alpha_{1}=\left(x_{1}, x_{2}, \cdots, x_{n}\right)^{\prime}, \beta=\left(y_{1}, y_{2}, \cdots, y_{n}\right)^{\prime}\) ，其中 \(x_{1}, y_{1} \neq 0\) ，且 \(\alpha^{\prime} \beta=1\) ，求矩阵 \(A=\beta \alpha^{\prime}\) 的特征值和特征向量．

\subsection{题目5}
5．设二次型 \(f\left(x_{1}, x_{2}, x_{3}\right)=x_{1}^{2}+x_{2}^{2}+x_{3}^{2}+2 \alpha x_{1} x_{2}+2 x_{1} x_{3}+2 \beta x_{2} x_{3}\) 在正交线性替换 \(X=Q Y\) 下化为标准形 \(y_{1}^{2}+2 y_{2}^{2}\) ，求 \(\alpha, \beta\) 的值和正交矩阵 \(Q\) 。

\subsection{题目6}
6．给定 \(P^{3}\) 中的向量\\
\(\alpha_{1}=(1,-2,1)^{\prime}, \alpha_{2}=(0,1,-2)^{\prime}, \alpha_{3}=(1,1,2)^{\prime}, \beta_{1}=(0,1,1)^{\prime}, \beta_{2}=(1,-4,3)^{\prime}\) 。（数据可能有误）记 \(W_{1}=L\left(\alpha_{1}, \alpha_{2}, \alpha_{3}\right), W_{2}=L\left(\beta_{1}, \beta_{2}\right)\) ，求 \(W_{1}+W_{2}\) 与 \(W_{1} \cap W_{2}\) 的基与维数。

\subsection{题目7}
7．已知 \(\sigma\) 是数域 \(P\) 上线性空间 \(V\) 的线性变换，\(h(x) \in P[x]\) 满足 \(h(x)=f(x) g(x),(f(x), g(x))=1\)，证明：

\[
\operatorname{Ker} h(\sigma)=\operatorname{Ker} f(\sigma) \oplus \operatorname{Ker} g(\sigma)
\]
\end{document}
